\documentclass[10pt]{article}

\usepackage[utf8]{inputenc}

\usepackage[T1]{fontenc}

\usepackage{amsmath}

\usepackage{amsfonts}

\usepackage{amssymb}

\usepackage[version=4]{mhchem}

\usepackage{stmaryrd}

\usepackage{bbold}

\usepackage{mathrsfs}



\begin{document}

При вещественных $k$



\[

\operatorname{Re} \sigma^{ \pm}(k, x)=-\frac{1}{2} \ln \left(1+\frac{\operatorname{Im} \sigma^{ \pm}(k, x)}{k}\right) .

\]



Если потенциал $q(x)=u(x, 0)$ принадлежит пространству Шварца (то есть бесконечно дифференцируем и при $|x| \rightarrow \infty$ убывает вместе со всеми производными быстрее любой степени $|x|^{-1}$ ), то решение $u(x, t)$ задачи (1) при каждом $t \in(-\infty, \infty)$ тоже принадлежит этому пространству, а функции $\sigma^{ \pm}(k, x)$ удовлетворяют одному и тому же асимптотич. равенству



\[

\sigma^{ \pm}(k, x) \simeq \sum_{p=1}^{\infty} \frac{\sigma_{p}(x, t)}{(2 i k)^{p}}, \quad|k| \rightarrow \infty, \quad \pm \operatorname{Im} k \geqslant 0,

\]



в к-ром коэффициенты $\sigma_{p}=\sigma_{p}(x, t)$ являются многочленами от $u, u_{x}, \ldots, u_{x \ldots x}^{(p-1)}$ и принадлежат пространству Шварца. Функции $\sigma_{p}(x, t)$ с четными номерами $p=2 q$ являются производными по $x$ от функций, принадлежащих пространству Шварца. Поэтому



\[

\ln a(k) \simeq-2 \sum_{q=0}^{\infty} \frac{I_{q-1}(u)}{(2 i k)^{2 q+1}}

\]



где



\[

I_{q-1}(u)=\frac{1}{2} S_{2 q+1}=\frac{1}{2} \int_{-\infty}^{\infty} \sigma_{2 q+1}(x, t) d x,

\]



причем это асимптотич. равенство выполняется независимо от того, удовлетворяет функция $u(x, t)$ уравнению Кортевега - де Фриса или нет.



Если $u(x, t)$ - решение задачи (1), то $\ln a(k)$, а значит, и коэффициенты $I_{q-1}(u)$ не зависят от $t$, то есть функционалы $I_{q-1}(u)$ образуют счетное семейство нетривиальных интегралов движения. Они называются полиноми альным и интегралам и д и жения, так қак являются интегралами от полиномов относительно функции $u(x, t)$ и ее производных по $x$. Напр.,



\[

\left.\begin{array}{c}

I_{-1}(u)=\frac{1}{2} \int_{-\infty}^{\infty} u(x, t) d x, \quad I_{0}(u)=-\frac{1}{2} \int_{-\infty}^{\infty} u(x, t)^{2} d x, \\

I_{1}(u)=\frac{1}{2} \int_{-\infty}^{\infty}\left(u_{x}(x, t)^{2}+2 u(x, t)^{3}\right) d x

\end{array}\right\}

\]



Их можно также выразить через корни $i x_{1}, i x_{2}, \ldots, i x_{m}$, $x_{j}>0$, функции $a(k)$ и значения ее модуля $|a(\xi)|=|a(-\xi)|$ на вещественной оси:



\[

\begin{aligned}

& I_{q-1}(u)=-\frac{2^{2 q+1}}{2 q+1} \sum_{p=1}^{m} x_{p}^{2 q+1}+ \\

& +2^{2 q+1} \frac{(-1)^{q}}{\pi} \int_{0}^{\infty} k^{2 q} \ln |a(k)| d k .

\end{aligned}

\]



В. Е. Захаров и Л.Д. Фаддеев [4] обнаружили глубокую связь между О. з. р.м. и теорией гамильтоновых систем, установив, что переход от функций $u(x, t) \mathrm{K}$ данным рассеяния операторов (2) является переходом к переменным типа действие-угол в соответствуюшей гамильтоновой системе. Для уравнения Кортевега - де Фриса фазовым пространством служит пространство Шварца, снабженное кососкалярным произведением



\[

[u, v]=\int_{-\infty}^{\infty} u(x) \frac{d v(x)}{d x} d x=\frac{1}{2} \int_{-\infty}^{\infty}(u d v-v d u),

\]



превращаюшим его в симплектич. линейное пространство $M$. Переменная $x$ играет роль номера координаты, а значение $u(x)$ функции $u$ является $x$-й координатой точки $u \in M$. Роль функций на фазовом пространстве играют функционалы $F(u)\left(F: M \rightarrow \mathbb{R}^{1}\right)$, а роль частных производных - ва- риационные производные $\frac{\delta F}{\delta u(x)}$, то есть такие обычные или обобщенные функции, что



\[

F(u+\varepsilon)-F(u)=\int_{-\infty}^{\infty} \frac{\delta F}{\delta u(x)} \varepsilon(x) d x

\]



с точностью до членов более высокого порядка малости, чем $\varepsilon(x)$. Функционал $F$ называется гладки м, если его вариационная производная принадлежит пространству Шварца.



Пуассонова структура на $M$ задается скобками Пуассона



\[

\{F, \mathscr{g}\}=\int_{-\infty}^{\infty} \frac{\delta F}{\delta u(x)} \frac{d}{d x}\left(\frac{\delta \mathfrak{I}}{\delta u(x)}\right) d x,

\]



К-рые удовлетворяют обычным требованиям.



Гладкие функционалы $H$ порождают уравнения с частными производными



\[

\frac{d}{d t} u=\frac{d}{d x}\left(\frac{\delta H}{\delta u(x)}\right)

\]



к-рые называются гамильтоновыми уравнениям и д в и жен и я с гамильтонианом $H$. Напр., согласно (24), (24')



\[

\frac{\delta I_{-1}}{\delta u(x)}=\frac{1}{2}, \quad \frac{\delta I_{0}}{\delta u(x)}=-u(x), \quad \frac{\delta I_{1}}{\delta u(x)}=-u_{x x}+3 u^{2}

\]



и гамильтоновы уравнения с гамильтонианами $H=I_{p}$, $p=-1,0,1, \ldots$, таковы:



\[

u_{t}=0, u_{t}=-u_{x}, u_{t}=6 u u_{x}-u_{x x x} \text {. }

\]



Таким образом, уравнение Кортевега - де Фриса является гамильтоновым с гамильтонианом $H=I_{1}$.



Значения функщионалов $F$ на траекториях $u(x, t)$ гамильтоновых уравнений (26) удовлетворяют уравнениям



\[

\frac{d F}{d t}=\{F, H\} .

\]



При замене координат $u(x) \rightarrow F(k)$ на фазовом пространстве новые координаты являются функционалами от старых. Поэтому уравнение (26) в новых координатах переходит в уравнение



\[

\frac{d F(k)}{d t}=\{F(k), H\} .

\]



Лежащий в основе О. 3. p. м. переход от функций $u(x) \mathrm{k}$ данным рассеяния является нек-рой заменой координат на фазовом пространстве. Вместо данных рассеяния в качестве новых координат удобнее взять величины



\[

\begin{gathered}

n(k)=\frac{2 k}{\pi} \ln a(k) a(-k), \quad \varphi(k)=\frac{1}{2 i} \ln \frac{b(k)}{b(-k)}, \quad 0<k<\infty, \\

N(p)=x_{p}^{2}, \quad \Phi(p)=-\ln b_{p} \bar{b}_{p}, \quad p=1,2, \ldots, m,

\end{gathered}

\]



через к-рые легко выразить данные рассеяния.



Для того чтобы найти вид уравнения Кортевега - де Фриса в этих новых координатах, нужно вычислить скобки Пуассона между ними и гамильтонианом $H=I_{1}$, а так как, согласно (25), все функционалы $I_{p}$ выражаются только через $n(k), N(p)$, то достаточно вычислить всевозможные скобки Пуассона между новыми координатами:



\[

\left\{n(k), n\left(k^{\prime}\right)\right\},\left\{n(k), \varphi\left(k^{\prime}\right)\right\}, \ldots

\]



Решения Йоста находятся из уравнений



\[

e^{ \pm}(k, y)=e^{i k y}-\int_{ \pm \infty}^{y} \frac{\sin k(\xi-y)}{k} u(\xi) e^{ \pm}(k, \xi) d \xi

\]



и при фиксированных $k$ и $\boldsymbol{y}$ являются функционалами от потенциала $u(y)$; их вариационные производные:



\[

\begin{aligned}

& \frac{\delta e^{+}(k, y)}{\delta u(x)}=\left\{\begin{array}{l}

0, \quad x<y, \\

-s(k, y, x) e^{+}(k, x), \quad x>y,

\end{array}\right. \\

& \frac{\delta e^{-}(k, y)}{\delta u(x)}=\left\{\begin{array}{l}

s(k, y, x) e^{-(k, x)}, \quad x<y, \\

0, \quad x>y,

\end{array}\right.

\end{aligned}

\]





\end{document}